\documentclass[10pt,a4paper]{amsart}
\usepackage[margin=19mm]{geometry}
\usepackage{amsmath,amssymb,mathtools}
\usepackage[scr=rsfs,cal=euler]{mathalfa}
\usepackage{xfrac}
\usepackage[unicode,hidelinks,bookmarks=false]{hyperref}
\newcommand{\ie}{i.e.\ }
\newcommand{\cf}{cf.\ }
\newcommand{\resp}{respectively\ }
\newcommand{\Subsection}{\subsection}
\providecommand{\nobreakdash}{\nobreak\hbox{-}}
\setlength{\emergencystretch}{2em}
\multlinegap0pt
\usepackage{bm}
\usepackage{bbm}
\usepackage{dsfont}
\usepackage{xspace}
\usepackage{cancel}
\usepackage{mathtools}
\usepackage{amsthm}
\makeatletter
\@ifundefined{theorem}{\newtheorem{theorem}{Theorem}}{}
\makeatother
\makeatletter
\newcommand{\R}{\mathbb{R}}
\expandafter\ifx\csname cproof\endcsname\relax
\newenvironment{cproof}
{\begin{proof}
 [Proof of Claim.]
 \vspace{-1.2\parsep}}
{\renewcommand{\qed}{\hfill $\Diamond$} \end{proof}}
\fi
\expandafter\ifx\csname subproof\endcsname\relax
\newenvironment{subproof}
{\begin{proof}
 [Proof of Subclaim.]
 \vspace{-1.2\parsep}}
{\renewcommand{\qed}{\hfill $\bigtriangledown$} \end{proof}}
\fi
\makeatother
\title[A Polyhedral Perspective on the Perfect Matching Lattice]{A Polyhedral Perspective on the Perfect Matching Lattice}
\author{Olha Silina}
\date{Source: arXiv:2511.03863v1 (CC BY 4.0)}
\begin{document}
\maketitle

\noindent\emph{Source and licence.} A Polyhedral Perspective on the Perfect Matching Lattice. Version arXiv:2511.03863v1 (CC BY 4.0), distributed under the Creative Commons Attribution 4.0 International licence, as recorded in the arXiv licence metadata for this exact version and shown on the abstract page https://arxiv.org/abs/2511.03863v1. Full rights notice: licenses/arxiv-2511.03863-CC-BY-4.0.txt.\par\smallskip

\noindent\textit{Editorial scope.} Counted unit: the Theorem whose source label is \texttt{thm:facet-nonbip} in the article (source0.tex lines 552--558, source label \texttt{thm:facet-nonbip}), delivered together with the article's own complete original proof of it (source0.tex lines 559--577, one \texttt{proof} environment of 19 lines). No mathematics is written, repaired or re-derived: statement and proof are the article's own text line for line. Required context retained uncounted: eq:pm-bipartite (lines 166--168). Every definition, display and label of those lines is retained unchanged. Omitted from this package and not redistributed: 1--165, 169--551, 578--605. Nothing inside a retained excerpt is omitted or altered: every retained body is delivered exactly as the article's lines are written, so no omission receipt of any kind accompanies this package. Citations: the retained excerpts carry no citation command at all, so no bibliography entry of the article is transcribed and no reference is redistributed. Reference adaptation: none was needed and none was made; every reference inside the delivered unit resolves inside it, so no unresolved reference or citation is produced. Printed slips kept literal: no printed word, symbol, hypothesis or number of the counted statement or of its proof was altered. Layout and class: the article's own class is \texttt{article} with packages including geometry, amsmath, amssymb, amsthm, amsfonts, enumerate, times, epsfig, color, hyperref, tikz, cleveref, mathrsfs, algpseudocode; the wrapper is the lane's pinned \texttt{amsart} editorial wrapper at 10pt on A4, which restates the theorem-like environments the retained text needs and adds the article's own macro definitions that the retained text needs (\texttt{\textbackslash R}), with extra wrapper packages (bm, bbm, dsfont, xspace, cancel, mathtools). The delivered numbering is the unit's own. Assets: no retained span contains a figure environment, a graphics-inclusion command, a PDF-page inclusion, a table or a TikZ picture; no image file is copied into this package, the package declares no asset dependency and no image file is redistributed with it. Licence: version arXiv:2511.03863v1 of this article is distributed under the Creative Commons Attribution 4.0 International licence, as recorded in the arXiv licence metadata for this exact version and shown on the abstract page https://arxiv.org/abs/2511.03863v1. A full rights notice is delivered at licenses/arxiv-2511.03863-CC-BY-4.0.txt. Characters: every retained excerpt is pure ASCII, so the delivered engine, XeLaTeX with its default Latin Modern text font, reports no missing character.
The main idea of the paper is to utilize odd cut decomposition methods to reduce our graph to a number of graphs whose perfect matching polytope coincides with its \emph{bipartite relaxation}:\begin{equation}\label{eq:pm-bipartite}
    P(G):=\left\{x\in \R^A_{+}: x(\delta(v))= 1  \forall v\in V \right\}.
\end{equation} 

We begin with a useful characterization. \begin{theorem}\label{thm:facet-nonbip}
    Let $G$ be a BvN brick. If an inequality $x_e\geq 0$ is a facet defining inequality for $PM(G)$ then one of the following holds: \begin{itemize}
        \item[(a)] $G-e$ is perfect matching covered and is a near-brick;
        \item[(b)] there is a unique edge $e'\neq e$ such that $G-e-e'$ is bipartite and $x_e = x_{e'}$ for all $x \in PM(G)$.
    \end{itemize}
    Furthermore, if $G$ is a BvN brick with an edge $e$ such that (a) or (b) holds, then $x_e\geq 0$ is a facet defining inequality for PM(G).
\end{theorem}

\begin{proof}
    We proceed by computing the affine dimension of $PM(G)$ and its facet. Notice that the equalities $\delta(v)=1$ are either independent or have rank $|V|-1$ if $G$ is bipartite. Thus, the dimension of (\ref{eq:pm-bipartite}) is exactly $|E|-|V|+\sigma_b(G)=|E|-|V|$, where $\sigma_b(H)$ denotes the number of bipartite connected components of $H$. 
    Now, consider a face defined by setting $x_e=0$. All vertices of this face are exactly the perfect matchings that do not use $e$ or, in other words, the perfect matchings of $G-e$. If $G-e$ is not perfect matching covered, then there is at least one edge $e'$ which is no longer in any perfect matching. Let $E'$ be the set of edges that belong to some perfect matching of $G-e$. Then, the vertices of the face defined by $x_e=0$ are exactly the perfect matchings of $G[E']$, so $PM(G[E'])$ has dimension $|E'|-|V|+\sigma_b(G[E'])$. Since this is one smaller than the original dimension, we get the following \begin{equation*}
        \sigma_b(G[E'])=|E\backslash E'|-1.
    \end{equation*} Let us denote this quantity by $k$. This means that by removing $k+1$ edges from $G$, the graph has at least $k$ bipartite components left. We will prove that unless $k\geq 2$ there is a $2$-separation in $G$, contradicting the fact that $G$ was a brick.

    Indeed, let $C_1, C_2, \ldots, C_k$ be the bipartite components. First, suppose $|\delta(C_i)|\leq 2$ for some component. It cannot be less than $2$ since it would mean that in $G$, there is a cut-edge, so it does not belong to any perfect matching. Then, let $u\in C_i$ and $v\notin C_i$ be the endpoints of two different edges leaving $C_i$. It is easy to see that $G-u-v$ has at least two odd components, so $\{u,v\}$ forms a $2$-separation. Now, we have $\delta(C_i)\geq 3$. Counting this over all components, we have \begin{equation}\label{eq:components-count}
        3k\leq \sum_i |\delta(C_i)|=2|E\backslash E'|-N \leq 2(k+1),
    \end{equation} where $N=0$ if $C_i$ are the only components and $N=\sum_j \delta(D_j)$ if $G[E']$ has more non-bipartite components $D_j$ (notice that there are at most two of those). From (\ref{eq:components-count}) we deduce $k\leq 2$.

    \underline{$k=2:$} in this case, (\ref{eq:components-count}) holds with equality, so there are no non-bipartite components. Let $A_i,B_i$ be the bipartition of $C_i$ for $i=1,2$. There are three edges $e,e_1,e_2$ in $E\backslash E'$ and $x_e=0$ implies $x_{e_1}=x_{e_2}=0$ for all $x \in PM(G)$.
    Without loss of generality, let $e$ connect $A_1$ to $B_2$. Then both $e_1$ and $e_2$ must connect $B_1$ to $A_2$, as otherwise there cannot be a perfect matching using $e$ and one of $e_1, e_2$. However, this makes the original graph bipartite, which is not possible.
    
    \underline{$k=1$:} this means that $E'=E\backslash\{e, e_1\}$ and $G-e-e_1$ is a perfect matching covered bipartite graph. Moreover, $e$ and $e_1$ are inside the opposite sides of the bipartition.

    \underline{$k=0$:} this means $G-e$ is perfect matching covered.

    To prove the converse, it suffices to redo the dimension calculations from the above part.
\end{proof}

\end{document}
